% ===========================================================================
% Section II-C -- Hilbert Demodulation and Windowed SES
% Owner: Member 3
% ===========================================================================
\subsection{Hilbert Demodulation and Windowed Squared Envelope Spectrum}
The bandpass stage leaves $y_{\text{BP}}[n]$ as a series of short bursts of
structural ringing. The bursts repeat at the fault rate $f_{\text{fault}}$.
The FFT of $y_{\text{BP}}[n]$ puts the energy around the resonance $\omega_r$,
so the fault rate is almost invisible in the spectrum. It has to be moved
down to a low-frequency region, where the harmonics are far apart and easy to
separate. That operation is demodulation.

The first step is the analytic signal, built with the Discrete Hilbert
Transform:
\begin{equation}
  z[n] = y_{\text{BP}}[n] + j\hat{y}_{\text{BP}}[n],
  \label{eq:analytic}
\end{equation}
where $\hat{y}_{\text{BP}}[n]=\mathcal{H}\{y_{\text{BP}}[n]\}$ is the Hilbert
transform of the filtered signal. A real record has a spectrum on both sides
of $0$~Hz, and the two halves mirror each other. The transform keeps the
positive half and drops the rest, which is what makes the envelope
well defined. It works as a $90^\circ$ phase shifter,
\begin{equation}
  H_{\text{Hilbert}}(e^{j\omega}) = -j\,\operatorname{sgn}(\omega),
  \label{eq:hilbert}
\end{equation}
so every positive frequency is shifted by $-\pi/2$ and the negative
frequencies are removed.

The usual envelope takes a square root of $y_{\text{BP}}^2+\hat{y}_{\text{BP}}^2$.
The squared envelope is used here instead:
\begin{equation}
  e[n] = |z[n]|^2 = z[n]z^{*}[n] = y_{\text{BP}}^2[n] + \hat{y}_{\text{BP}}^2[n].
  \label{eq:se}
\end{equation}
Squaring does two useful things. A square in time is an
auto-convolution in frequency, so the modulation theorem moves the carrier
$\omega_r$ down to baseband ($0$~Hz) without distortion. A square root does
not behave this way. It is non-linear, and the intermodulation harmonics it
adds can hide a weak fault. Squaring also favours large amplitudes: an impact
peak is squared into a much larger number, while the background changes
little. The fault-to-noise ratio therefore improves before any threshold is
applied.

The squared envelope is never negative, so it carries a large DC term. That
term is much bigger than the fault harmonics, and its mainlobe leaks into the
low-frequency bins where the cage and shaft harmonics sit. The mean is removed
first:
\begin{equation}
  e_0[n] = e[n] - \mu_e, \qquad
  \mu_e = \frac{1}{N}\sum_{n=0}^{N-1}e[n].
  \label{eq:emean}
\end{equation}

A fault frequency rarely sits exactly on a DFT bin centre. At $107.36$~Hz,
with $\Delta f=f_s/N_{\text{FFT}}$, the mismatch is unavoidable. Rectangular
truncation then spreads sidelobes into the neighbouring bins. An $N$-point
Hann window controls them:
\begin{equation}
  w[n] = \tfrac12\left(1-\cos\frac{2\pi n}{N-1}\right),
  \qquad n = 0,1,\dots,N-1.
  \label{eq:hann}
\end{equation}
The window brings both ends to zero, which halves the coherent amplitude. It
also widens the mainlobe, so a little frequency resolution is lost, and that
trade is accepted for the much lower sidelobes. The coherent gain
\begin{equation}
  CG = \frac{1}{N}\sum_{n=0}^{N-1}w[n] = \tfrac12
  \label{eq:cg}
\end{equation}
puts the amplitude back, so the SES levels keep their physical meaning. The
one-sided SES is then computed with an
$N_{\text{FFT}}$-point radix-2 FFT:
\begin{equation}
  \text{SES}[k] = \frac{2}{N\cdot CG}\left|\sum_{n=0}^{N-1} w[n]\,e_0[n]\,
  e^{-j\frac{2\pi}{N_{\text{FFT}}}kn}\right|,
  \label{eq:ses}
\end{equation}
for $k=0,1,\dots,N_{\text{FFT}}/2$, where bin $k$ corresponds to
$f_k=k\,f_s/N_{\text{FFT}}$~Hz.
